\chapter{Thống kê và chất lượng đồ thị}
\label{ch:quality}
\section{Thành phần đồ thị chính thức}
V4 chính thức gồm \TotalNodes{} nút và \TotalEdges{} cạnh được chấp nhận.
Bảng~\ref{tab:nodes} và~\ref{tab:edges} trình bày thống kê theo loại. Trong
20 chủ đề có bốn chủ đề cha và 16 chủ đề lá; hai nút hội nghị biểu diễn chuỗi
hội nghị, còn năm được lưu trên nút bài báo.

% Generated from frozen v4; edit generator.
\begin{table}[!htbp]\centering\small
\caption{\lang{Graph nodes by type.}{Số nút theo loại thực thể.}}\label{tab:nodes}
\begin{tabularx}{\textwidth}{Yr}
\toprule \lang{Entity type}{Loại thực thể} & \lang{Nodes}{Số nút}\\\midrule
\lang{Organization}{Tổ chức} & \num{2047}\\
\lang{Paper}{Bài báo} & \num{13617}\\
\lang{Researcher}{Nhà nghiên cứu} & \num{32189}\\
\lang{Topic}{Chủ đề} & \num{20}\\
\lang{Venue}{Hội nghị} & \num{2}\\
\textbf{\lang{Total}{Tổng}} & \num{47875}\\
\bottomrule\end{tabularx}\end{table}

% Generated from frozen v4; edit generator.
\begin{table}[!htbp]\centering\small
\caption{\lang{Accepted edges by predicate and evidence class.}{Các cạnh được chấp nhận theo quan hệ và loại bằng chứng.}}\label{tab:edges}
\begin{tabularx}{\textwidth}{Ylr}
\toprule \lang{Predicate}{Quan hệ} & \lang{Evidence}{Bằng chứng} & \lang{Edges}{Số cạnh}\\\midrule
\texttt{AUTHORED} & \lang{Observed}{Quan sát} & \num{50437}\\
\texttt{AFFILIATED\_WITH} & \lang{Observed}{Quan sát} & \num{38664}\\
\texttt{PUBLISHED\_IN} & \lang{Observed}{Quan sát} & \num{13617}\\
\texttt{CITES} & \lang{Observed}{Quan sát} & \num{7308}\\
\texttt{HAS\_TOPIC} & \lang{Inferred}{Suy ra từ mô hình} & \num{14222}\\
\texttt{SUBTOPIC\_OF} & \lang{Curated}{Biên soạn} & \num{16}\\
\texttt{COAUTHORED\_WITH} & \lang{Derived}{Dẫn xuất} & \num{80482}\\
\texttt{COLLABORATED\_WITH} & \lang{Derived}{Dẫn xuất} & \num{13031}\\
\texttt{RESEARCHES\_TOPIC} & \lang{Derived}{Dẫn xuất} & \num{5622}\\
\texttt{CONTRIBUTES\_TO\_TOPIC} & \lang{Derived}{Dẫn xuất} & \num{2332}\\
\textbf{\lang{Total}{Tổng}} &  & \num{225731}\\
\bottomrule\end{tabularx}\end{table}

\FloatBarrier

Nhóm quan hệ thư mục quan sát được có 110,026 cạnh. Các nhóm còn lại gồm
14,222 cạnh chủ đề suy ra, 16 cạnh phân cấp biên soạn và 101,467 cạnh dẫn xuất.
Đặc biệt, số quan hệ tác giả, số liên kết tổ chức tổng hợp và số cặp đồng tác
giả đo những đối tượng khác nhau. Không thể cộng chúng rồi diễn giải thành số
quan sát khoa học độc lập.

\Needspace{15cm}
\section{Kiểm tra xây dựng và chất lượng bằng chứng}
Bộ kiểm tra cấu trúc đã ghi nhận kết quả đạt cho v4: ID duy nhất, nhãn và quan
hệ hợp lệ, loại thực thể ở hai đầu đúng quy định, trạng thái chấp nhận, tham
chiếu bằng chứng và phép tính lại số bài hỗ trợ/trọng số của quan hệ dẫn xuất.
Cả mười quan hệ đều hiện diện. Lần nạp Neo4j được ghi nhận đã đối soát số nút
và cạnh; sáu truy vấn toàn vẹn và phép nạp lại không làm thay đổi dữ liệu cũng
đạt. Đây là hồ sơ kiểm tra của phiên bản, không phải các lần kiểm tra cơ sở dữ
liệu mới thực hiện khi viết báo cáo.

Kiểm tra tính nhất quán toàn bộ nguồn đã so sánh đối tượng tác giả và tham
chiếu được giữ lại cho cả 13,606 công trình đối sánh, bao phủ 65,469 vị trí và
341 mã băm phản hồi lưu đệm. Kết quả xác nhận sự trùng khớp với ảnh chụp nguồn
đã lưu. Tính nhất quán cấu trúc và byte nguồn không chứng minh định danh hay
đơn vị công tác do nhà cung cấp ghi nhận là đúng.

Đợt rà soát tự động dựa trên nguồn, được thiết kế trước khi đánh giá, nhận định
193 trong 200 cạnh được lấy mẫu ngẫu nhiên đều là có hỗ trợ (96.50\%; khoảng
Wilson 95\%: 92.95--98.29\%), với một trường hợp được đánh giá là sai và sáu
trường hợp chưa thể xác minh.
Trường hợp chưa xác minh được được tính là không có hỗ trợ. Đây là
\textbf{đánh giá có hỗ trợ AI}, không phải nhãn do người gán hay kiểm chứng độc
lập: nguồn dùng khi xây dựng được sử dụng lại, tính độc lập của bên đánh giá
chưa được xác lập, và khoảng thống kê chỉ phản ánh biến động lấy mẫu chứ không
bao gồm sai số đánh giá. Một chẩn đoán phân tầng theo quan hệ thực hiện sau đó
đánh dấu ba liên kết dẫn xuất từ chủ đề và một quan hệ tác giả chưa chắc chắn;
nó không sửa đồ thị cố định và không được gộp vào ước lượng mẫu ngẫu nhiên đều.

Các khoảng trống ở Bảng~\ref{tab:coverage} ảnh hưởng đến phần mạng quan sát
được. Có 199 bài không có quan hệ tác giả được chấp nhận, và bằng chứng tổ chức
còn thiếu ở một phần danh sách tác giả chính thức. Độ bao phủ chủ đề cũng đo
khả năng có nhãn, không đo tính đúng của nhãn. Trên 50 bài kiểm tra lịch sử,
nhãn dùng vector có độ chính xác dự đoán dương (precision) vi mô 0.6863, độ thu hồi (recall) vi mô 0.3535 và F1 0.4667 khi so với nhãn tham chiếu do AI gán; F1 của phương
pháp cơ sở dùng cụm từ là 0.2301. Các nhãn tham chiếu đã xuất hiện trong các lần rà soát trước này không tạo
thành một tập kiểm tra chủ đề độc lập do con người thực hiện.
Vì vậy, xu hướng và hồ sơ chủ đề mang tính thăm dò; số bài quan sát và thống kê
cạnh được chấp nhận vẫn là các đại lượng đo trực tiếp được.
